% !TEX root = master.tex
\chapter{Verzeichnisstruktur (Backend / Webserver)}
\label{ch:a-verzeichnis}

Das Projektverzeichnis wird vollst\"andig in das Webserver-Wurzelverzeichnis (bei XAMPP
\texttt{htdocs}) installiert. Alle internen Pfade sind relativ aufgel\"ost, sodass das Projekt
unabh\"angig vom Ordnernamen lauff\"ahig ist. Die folgende \"Ubersicht zeigt die Struktur:

\begin{lstlisting}[numbers=none]
guardx/
|-- index.php                 Startseite / Dashboard
|-- init.php                  Session-Start, Timeout, CSRF-Token
|-- config.php                Zentrale Konfiguration (DB-Verbindung)
|-- functions.php             Server-Hilfsfunktionen (IP, Pfade, Mobile)
|-- functions.js              Client-Logik (Session-Timer, Menue)
|-- auth.php                  Login / Registrierung / Logout
|-- user.php                  Persoenlicher Bereich (Profil)
|-- admin.php                 Admin-Bereich (Benutzerverwaltung, CRUD)
|-- admin_logout.php          Logout-Endpunkt
|-- keep_alive.php            Session am Leben halten (AJAX-Ziel)
|-- impressum.php             Impressum
|-- datenschutz.php           Datenschutzerklaerung
|-- agb.php                   AGB
|-- php.ini.hinweise          Hinweise zu benoetigten PHP-Extensions
|
|-- db/
|   |-- users.sql             Aktueller DB-Dump (Struktur + Daten)
|   |-- init_databases.sql    Aelterer Dump (siehe Hinweis Kap. 8)
|
|-- include/
|   |-- api-calls/
|   |   |-- skript.php         Oberflaeche Live-Passwort-Checker
|   |   |-- check.php          AJAX-Endpunkt: Staerke + Leak-Abgleich
|   |-- fingerprinting/
|   |   |-- info.php           Modul Browser-Fingerprinting
|   |-- metadata_stripping/
|       |-- metadata_stripping.php  Modul Metadaten-Bereiniger
|       |-- uploads/           (zur Laufzeit erzeugt) Upload-Ablage
|       |-- clean/             (zur Laufzeit erzeugt) bereinigte Bilder + JSON
|
|-- template/
|   |-- navbar.php             Desktop-Navigation
|   |-- navbar_mobile.php      Mobile Navigation (Hamburger)
|   |-- footer.php             Footer mit Rechts-Links
|
|-- style/
    |-- main.css               Basis-/Layout-Styles (Cyber-Theme)
    |-- navbar.css             Navigationsstyles
    |-- metadata_stripping.css Styles des Metadaten-Moduls
\end{lstlisting}

\noindent Die Verzeichnisse \texttt{include/metadata\_stripping/uploads/} und
\texttt{.../clean/} werden vom Modul bei Bedarf automatisch angelegt; sie m\"ussen f\"ur den
Webserver-Benutzer beschreibbar sein. Tabelle~\ref{tab:a-verzeichnis} fasst die Aufgaben der
wichtigsten Komponenten zusammen.

\begin{small}
\begin{longtable}{@{}L{4.4cm} L{9.9cm}@{}}
\caption{Komponenten von Projekt~A}
\label{tab:a-verzeichnis}\\
\toprule
\textbf{Datei / Verzeichnis} & \textbf{Aufgabe} \\
\midrule
\endfirsthead
\toprule
\textbf{Datei / Verzeichnis} & \textbf{Aufgabe} \\
\midrule
\endhead
\bottomrule
\endfoot
\bottomrule
\endlastfoot
\texttt{init.php} & Startet/erneuert die Session, erzwingt das Inaktivit\"ats-Timeout und erzeugt das CSRF-Token. Wird von praktisch allen Seiten eingebunden. \\
\texttt{config.php} & Zentrale Konfigurationsdatei mit den Verbindungsdaten zur Datenbank; wird per \texttt{require\_once} eingebunden. \\
\texttt{functions.php} & Serverseitige Hilfsfunktionen: sichere IP-Ermittlung, Pfad-Helfer (\texttt{get\_url}), Mobile-Erkennung (\texttt{is\_mobile}). \\
\texttt{functions.js} & Clientseitige Logik: Session-Timer mit jQuery-UI-Dialog, mobiles Men\"u, Login-Weiterleitung. \\
\texttt{auth.php} & Registrierung, Login und Logout inkl.\ Validierung, Hashing, Brute-Force- und Rate-Limit-Schutz. \\
\texttt{user.php} & Pers\"onlicher Bereich: E-Mail/Passwort \"andern, Account l\"oschen. \\
\texttt{admin.php} & Admin-Bereich mit vollst\"andiger Benutzerverwaltung (Anlegen, \"Andern, Rolle, L\"oschen). \\
\texttt{include/api-calls/} & Live-Passwort-Checker: Oberfl\"ache (\texttt{skript.php}) und AJAX-Endpunkt (\texttt{check.php}). \\
\texttt{include/fingerprinting/} & Modul Browser-Fingerprinting (\texttt{info.php}). \\
\texttt{metadata\_stripping/} & Modul Metadaten-Bereiniger inkl.\ Lauf\-zeit-Ordnern f\"ur Upload und bereinigte Ausgabe. \\
\texttt{template/} & Wiederverwendbare Bausteine: Desktop-/Mobil-Navigation und Footer. \\
\texttt{style/} & Stylesheets f\"ur Layout, Navigation und Metadaten-Modul. \\
\texttt{db/} & SQL-Dumps der Datenbank (Struktur und Daten). \\
\bottomrule
\end{longtable}
\end{small}
